% Options for packages loaded elsewhere
\PassOptionsToPackage{unicode}{hyperref}
\PassOptionsToPackage{hyphens}{url}
\PassOptionsToPackage{dvipsnames,svgnames,x11names}{xcolor}
\documentclass[
  english,
  10pt,
  a4paper,
]{article}
\usepackage{xcolor}
\usepackage[top=22mm,bottom=20mm,left=18mm,right=18mm]{geometry}
\usepackage{amsmath,amssymb}
\setcounter{secnumdepth}{-\maxdimen} % remove section numbering
\usepackage{iftex}
\ifPDFTeX
  \usepackage[T1]{fontenc}
  \usepackage[utf8]{inputenc}
  \usepackage{textcomp} % provide euro and other symbols
\else % if luatex or xetex
  \usepackage{unicode-math} % this also loads fontspec
  \defaultfontfeatures{Scale=MatchLowercase}
  \defaultfontfeatures[\rmfamily]{Ligatures=TeX,Scale=1}
\fi
\usepackage{lmodern}
\ifPDFTeX\else
  % xetex/luatex font selection
\fi
% Use upquote if available, for straight quotes in verbatim environments
\IfFileExists{upquote.sty}{\usepackage{upquote}}{}
\IfFileExists{microtype.sty}{% use microtype if available
  \usepackage[]{microtype}
  \UseMicrotypeSet[protrusion]{basicmath} % disable protrusion for tt fonts
}{}
\makeatletter
\@ifundefined{KOMAClassName}{% if non-KOMA class
  \IfFileExists{parskip.sty}{%
    \usepackage{parskip}
  }{% else
    \setlength{\parindent}{0pt}
    \setlength{\parskip}{6pt plus 2pt minus 1pt}}
}{% if KOMA class
  \KOMAoptions{parskip=half}}
\makeatother
\usepackage{longtable,booktabs,array}
\usepackage{caption}
\captionsetup[table]{skip=6pt}
\newcounter{none} % for unnumbered tables
\usepackage{calc} % for calculating minipage widths
% Correct order of tables after \paragraph or \subparagraph
\usepackage{etoolbox}
\makeatletter
\patchcmd\longtable{\par}{\if@noskipsec\mbox{}\fi\par}{}{}
\makeatother
\ifLuaTeX
\usepackage[bidi=basic,shorthands=off]{babel}
\else
\usepackage[bidi=default,shorthands=off]{babel}
\fi
\ifLuaTeX
  \usepackage{selnolig} % disable illegal ligatures
\fi
\setlength{\emergencystretch}{3em} % prevent overfull lines
\providecommand{\tightlist}{%
  \setlength{\itemsep}{0pt}\setlength{\parskip}{0pt}}
% ---- Çılgın Kurye document theme: same look as the CSE 396 project proposal ----
% Compile with XeLaTeX (also works on Overleaf). Fonts: TeX Gyre Termes (Times), DejaVu Sans Mono.
\usepackage{fontspec}
\setmainfont{texgyretermes}[Extension=.otf, UprightFont=*-regular, BoldFont=*-bold, ItalicFont=*-italic, BoldItalicFont=*-bolditalic]
\setsansfont{texgyreheros}[Extension=.otf, UprightFont=*-regular, BoldFont=*-bold, ItalicFont=*-italic, BoldItalicFont=*-bolditalic]
\setmonofont{DejaVuSansMono}[Extension=.ttf, UprightFont=*, BoldFont=*-Bold, ItalicFont=*-Oblique, BoldItalicFont=*-BoldOblique, Scale=0.80]
\setmathfont{texgyretermes-math.otf}

\usepackage{xcolor}
\definecolor{ckblue}{HTML}{1F3A5F}
\definecolor{cklight}{HTML}{E8EEF5}
\definecolor{ckgrey}{HTML}{F2F2F2}
\definecolor{ckline}{HTML}{555555}
\definecolor{shadecolor}{HTML}{F4F5F7}

% Headings as in the proposal: serif, bold; subsections bold italic. Numbers are part of the heading text.
\setcounter{secnumdepth}{0}
\usepackage{titlesec}
\titleformat{\section}{\rmfamily\Large\bfseries}{}{0pt}{}
\titleformat{\subsection}{\rmfamily\large\bfseries\itshape}{}{0pt}{}
\titleformat{\subsubsection}{\rmfamily\normalsize\bfseries}{}{0pt}{}
\titlespacing*{\section}{0pt}{12pt}{4pt}
\titlespacing*{\subsection}{0pt}{8pt}{3pt}

% Header and footer (the first page has only the footer, like the proposal)
\usepackage{fancyhdr,lastpage}
\pagestyle{fancy}
\fancyhf{}
\fancyhead[L]{\footnotesize\itshape CSE 396 · Group 8 · Çılgın Kurye}
\fancyhead[R]{\footnotesize\itshape\ckmodule}
\fancyfoot[C]{\footnotesize Page \thepage\ of \pageref*{LastPage}}
\renewcommand{\headrulewidth}{0.4pt}
\setlength{\headheight}{14pt}
\fancypagestyle{ckfirst}{\fancyhf{}\renewcommand{\headrulewidth}{0pt}\fancyfoot[C]{\footnotesize Page \thepage\ of \pageref*{LastPage}}}

% Inline code: breaks only after . : ( , / _ and spaces, never inside a name
\usepackage[obeyspaces,spaces]{url}
\AtBeginDocument{%
  \renewcommand{\UrlBreaks}{\do\.\do\:\do\(\do\,\do\/\do\_\do\|}%
  \renewcommand{\UrlBigBreaks}{}%
  \renewcommand{\UrlNoBreaks}{}%
  \urlstyle{tt}}

% Tables: boxed grid as in the proposal, smaller text
\usepackage{etoolbox}
\AtBeginEnvironment{longtable}{\footnotesize\renewcommand{\arraystretch}{1.2}}
\setlength{\LTpre}{6pt}
\setlength{\LTpost}{8pt}
\setlength{\arrayrulewidth}{0.4pt}

% Code blocks: long lines wrap; diagrams stay in one piece
\usepackage{framed}
\usepackage{fvextra}
\RecustomVerbatimEnvironment{verbatim}{Verbatim}{fontsize=\footnotesize, breaklines}
\makeatletter
\@ifundefined{Shaded}{}{\renewenvironment{Shaded}{\begin{snugshade}}{\end{snugshade}}}
\@ifundefined{Highlighting}{}{\RecustomVerbatimEnvironment{Highlighting}{Verbatim}{commandchars=\\\{\}, fontsize=\footnotesize, breaklines, breakanywhere}}
\makeatother

% Figures drawn with TikZ
\usepackage{tikz}
\usetikzlibrary{arrows.meta,positioning,fit,calc,backgrounds,shapes.geometric}
\tikzset{
  ck/.style={font=\footnotesize, >={Stealth[length=5pt]}, line width=0.6pt},
  box/.style={draw=ckline, fill=white, rounded corners=2pt, align=center, inner sep=4pt, minimum height=7mm},
  mod/.style={box, fill=ckgrey},
  this/.style={box, fill=cklight, draw=ckblue, line width=0.9pt},
  group/.style={draw=ckline, dashed, rounded corners=3pt, inner sep=6pt},
  lbl/.style={font=\scriptsize, align=center, inner sep=1.5pt},
  msg/.style={->, draw=ckline},
}
\usepackage{caption}
\captionsetup[figure]{font=small, labelfont=bf, skip=4pt}
\newcommand{\ckfigure}[2]{\begin{figure}[htbp]\centering #1\caption{#2}\end{figure}}

\usepackage{enumitem}
\setlist{itemsep=1pt, topsep=3pt}
\setlength{\parskip}{4pt}
\setlength{\parindent}{0pt}
\raggedbottom
% ---- end of theme ----
\newcommand{\ckmodule}{M2 Network and Synchronization}
\usepackage{bookmark}
\IfFileExists{xurl.sty}{\usepackage{xurl}}{} % add URL line breaks if available
\urlstyle{same}
% fallback for those not using the hyperref driver hyperxmp:
\makeatletter
\@ifundefined{xmpquote}{\newcommand{\xmpquote}[1]{#1}}{}
\makeatother
\hypersetup{
  pdftitle={M2 Network and Synchronization},
  pdflang={en},
  colorlinks=true,
  linkcolor={black},
  filecolor={Maroon},
  citecolor={Blue},
  urlcolor={black},
  pdfcreator={LaTeX}}

\title{M2 Network and Synchronization --- Module Documentation}
\author{}
\date{}


% ---- colours and drawings ----
% ---- colours ----
\definecolor{srv}{HTML}{2E6FB7}\colorlet{srvL}{srv!15}
\definecolor{cli}{HTML}{2E9E5B}\colorlet{cliL}{cli!15}
\definecolor{adm}{HTML}{E08A1E}\colorlet{admL}{adm!18}
\definecolor{brd}{HTML}{7B4FB3}\colorlet{brdL}{brd!15}
\definecolor{tcpc}{HTML}{1F5FA8}
\definecolor{udpc}{HTML}{D9531E}
\definecolor{obs}{HTML}{D7263D}
\definecolor{tA}{HTML}{00A6D6}
\definecolor{tB}{HTML}{3BB273}
\definecolor{tC}{HTML}{F2C14E}
\definecolor{tD}{HTML}{9B5DE5}

% ---- styles ----
\tikzset{
  nb/.style={rounded corners=3pt, line width=0.9pt, align=center, inner sep=4pt, minimum height=9mm, font=\footnotesize},
  sv/.style={nb, draw=srv, fill=srvL},
  cl/.style={nb, draw=cli, fill=cliL},
  ad/.style={nb, draw=adm, fill=admL},
  bd/.style={nb, draw=brd, fill=brdL},
  hwb/.style={nb, draw=black!50, fill=black!5},
  tcpa/.style={draw=tcpc, line width=1.5pt},
  udpa/.style={draw=udpc, line width=1.5pt, dash pattern=on 5pt off 2.5pt},
  inna/.style={draw=black!55, line width=1pt},
  al/.style={font=\scriptsize, fill=white, inner sep=1.2pt, rounded corners=1pt},
  grp/.style={draw=black!35, dashed, rounded corners=5pt, inner sep=7pt},
  glbl/.style={font=\scriptsize\bfseries, text=black!65},
  hl/.style={line width=2.6pt},
  st/.style={nb, rounded corners=8pt, minimum width=24mm},
  co/.style={draw=black!60, line width=0.5pt},
  cot/.style={font=\scriptsize, align=left, text=black!80},
}

% ---- system map; argument = module to highlight (M1 ... M8, or none) ----
\newcommand{\systemmap}[1]{%
\tikzset{M1/.style={},M2/.style={},M3/.style={},M4/.style={},M5/.style={},M6/.style={},M7/.style={},M8/.style={},none/.style={}}%
\tikzset{#1/.style={hl}}%
\begin{tikzpicture}[ck]
\node[cl, M5, minimum width=27mm] (m5) at (0,3) {\textbf{M5}\\Player phones};
\node[ad, M6, minimum width=27mm] (m6) at (0,0.4) {\textbf{M6}\\Admin panel};
\node[sv, M2, minimum width=27mm, minimum height=38mm] (m2) at (4.6,1.7) {\textbf{M2}\\Network\\and sync};
\node[sv, M1, minimum width=27mm, minimum height=38mm] (m1) at (8.4,1.7) {\textbf{M1}\\Game server\\and rules};
\node[sv, M3, minimum width=25mm] (m3) at (12,2.8) {\textbf{M3}\\Physics};
\node[sv, M4, minimum width=25mm] (m4) at (12,0.6) {\textbf{M4}\\Traffic and AI};
\node[grp, fit=(m2)(m1)(m3)(m4)] (pc) {};
\node[glbl, anchor=south west] at (pc.north west) {MAIN COMPUTER};
\node[bd, M7, minimum width=27mm] (m7) at (4.6,-2.7) {\textbf{M7}\\ESP32 software};
\node[bd, M8, minimum width=27mm] (m8) at (8.4,-2.7) {\textbf{M8}\\LED engine};
\node[hwb, minimum width=25mm] (led) at (12,-2.7) {344 LEDs};
\node[hwb, minimum width=27mm] (sm) at (4.6,-4.4) {Sensor and\\4 motors};
\node[grp, fit=(m7)(m8)(led)(sm)] (bb) {};
\node[glbl, anchor=north east] at (bb.south east) {LED CITY BOARD};
\draw[tcpa, <->] (m5.east) -- node[al]{WebSocket} (m5.east -| m2.west);
\draw[tcpa, <->] (m6.east) -- node[al]{WebSocket} (m6.east -| m2.west);
\draw[tcpa, <->] ([xshift=-6mm]m7.north) -- node[al]{TCP} ([xshift=-6mm]m2.south);
\draw[udpa, <->] ([xshift=6mm]m7.north) -- node[al]{UDP} ([xshift=6mm]m2.south);
\draw[inna, <->] (m2) -- (m1);
\draw[inna, <->] (m1.east |- m3) -- (m3);
\draw[inna, <->] (m1.east |- m4) -- (m4);
\draw[inna, ->] (m7) -- node[al]{state} (m8);
\draw[inna, ->] (m8) -- (led);
\draw[inna, <->] (m7) -- (sm);
\draw[tcpa] (-1.3,-6.0) -- ++(0.9,0) node[right, font=\scriptsize]{TCP: every message arrives, in order};
\draw[udpa] (4.0,-6.0) -- ++(0.9,0) node[right, font=\scriptsize]{UDP: fast, the newest message wins};
\draw[inna] (9.3,-6.0) -- ++(0.9,0) node[right, font=\scriptsize]{inside one program or chip};
\end{tikzpicture}}

% ---- LED board seen from above (6 x 6 roads, 1 unit = 10 cm) ----
\newcommand{\boardbase}{%
\fill[black!88, rounded corners=3pt] (-0.5,-0.5) rectangle (5.5,5.5);
\foreach \i in {0,...,4}{\foreach \j in {0,...,4}{%
  \fill[black!60] (\i+0.13,\j+0.13) rectangle (\i+0.87,\j+0.87);
  \fill[black!22] (\i+0.24,\j+0.24) rectangle (\i+0.52,\j+0.52);
  \fill[black!30] (\i+0.56,\j+0.30) rectangle (\i+0.78,\j+0.62);
  \fill[black!38] (\i+0.28,\j+0.58) rectangle (\i+0.66,\j+0.78);}}
\foreach \i in {0,...,5}{%
  \draw[black, line width=4.5pt] (\i,-0.05) -- (\i,5.05);
  \draw[black, line width=4.5pt] (-0.05,\i) -- (5.05,\i);}
\foreach \i in {0,...,5}{\foreach \k in {0,...,30}{%
  \fill[black!55] (\i,\k/6) circle (0.025);
  \fill[black!55] (\k/6,\i) circle (0.025);}}
\foreach \x/\y/\c in {0/4/tA, 4/4/tB, 0/0/tC, 4/0/tD}{%
  \fill[\c!30!black!60] (\x+0.13,\y+0.13) rectangle (\x+0.87,\y+0.87);
  \fill[\c] (\x+0.3,\y+0.3) rectangle (\x+0.7,\y+0.7);
  \fill[white] (\x+0.4,\y+0.5) circle (0.05); \fill[white] (\x+0.6,\y+0.5) circle (0.05);}}
\newcommand{\nmark}[3]{\node[circle, fill=yellow!85!orange, draw=black, line width=0.4pt, inner sep=0.6pt, minimum size=3.6mm, font=\tiny\bfseries] at (#1,#2) {#3};}
\newcommand{\bike}[4]{%
\fill[#3!35] ($(#1,#2)+(#4:0.40)$) circle (0.05);
\fill[#3!65] ($(#1,#2)+(#4:0.21)$) circle (0.075);
\fill[#3, draw=white, line width=0.4pt] (#1,#2) circle (0.12);}
\newcommand{\dpoint}[3]{\draw[#3!50, line width=0.8pt] (#1,#2) circle (0.26); \draw[#3, line width=1.4pt] (#1,#2) circle (0.16); \fill[#3] (#1,#2) circle (0.07);}
\newcommand{\heavy}[2]{\draw[obs!60!black, line width=4pt] (#1) -- (#2); \draw[obs, line width=4pt, dash pattern=on 3pt off 5pt] (#1) -- (#2);}
\newcommand{\roadwork}[2]{\draw[obs, line width=4pt] (#1) -- (#2);}
\newcommand{\accident}[2]{\draw[obs, line width=4pt] (#1) -- (#2); \draw[white, line width=4pt, dash pattern=on 2.5pt off 3.5pt] (#1) -- (#2);}
\newcommand{\puddle}[2]{\draw[cyan!35, line width=4pt] (#1) -- (#2); \draw[white, line width=1.6pt, dash pattern=on 1pt off 2.5pt] (#1) -- (#2);}
\newcommand{\findme}[2]{\fill[white, opacity=0.25] (#1,#2) circle (0.36); \fill[white, opacity=0.5] (#1,#2) circle (0.24); \fill[white] (#1,#2) circle (0.13);}

% ---- phone game screen (7 x 3.4 units) ----
\newcommand{\phonemock}{%
\fill[black!85, rounded corners=9pt] (-0.3,-0.3) rectangle (7.3,3.7);
\fill[cyan!20] (0,1.7) rectangle (7,3.4);
\fill[cli!14] (0,0) rectangle (7,1.7);
\filldraw[fill=srv!30, draw=srv!70!black] (0.2,1.5) rectangle (1.0,3.0);
\filldraw[fill=adm!35, draw=adm!70!black] (1.1,1.6) rectangle (1.8,2.4);
\filldraw[fill=brd!30, draw=brd!70!black] (1.9,1.65) -- (2.6,1.65) -- (2.6,2.2) -- (2.25,2.5) -- (1.9,2.2) -- cycle;
\filldraw[fill=adm!35, draw=adm!70!black] (4.5,1.65) rectangle (5.1,2.3);
\filldraw[fill=srv!30, draw=srv!70!black] (5.2,1.55) rectangle (6.0,3.1);
\filldraw[fill=cli!35, draw=cli!70!black] (6.1,1.6) rectangle (6.85,2.6);
\fill[brown!60!black] (4.17,1.6) rectangle (4.23,1.8); \fill[cli!70!black] (4.2,1.95) circle (0.18);
\fill[black!55] (1.8,0) -- (5.2,0) -- (3.75,1.7) -- (3.25,1.7) -- cycle;
\draw[white, line width=0.8pt, dash pattern=on 4pt off 4pt] (3.5,0) -- (3.5,1.7);
\filldraw[fill=black!15, draw=black!70] (3.62,1.12) rectangle (3.88,1.32);
\fill[obs] (3.05,1.22) -- (3.15,1.48) -- (3.25,1.22) -- cycle;
\fill[black!85] (3.43,0.08) rectangle (3.57,0.3);
\fill[tA] (3.3,0.28) rectangle (3.7,0.6);
\fill[brown!65] (3.34,0.6) rectangle (3.66,0.82);
\fill[tA!60!black] (3.5,0.95) circle (0.13);
\node[fill=black!70, text=white, rounded corners=2pt, font=\scriptsize\bfseries, inner sep=2pt] at (3.5,3.15) {2:41};
\foreach \c/\s/\x in {tA/120/0.35,tB/95/0.95,tC/140/1.55,tD/80/2.15}
  \node[fill=\c, text=white, font=\tiny\bfseries, rounded corners=1pt, inner sep=1.5pt, minimum width=5mm] at (\x,3.15) {\s};
\fill[white, opacity=0.9] (6.45,2.85) circle (0.38); \draw[black!60] (6.45,2.85) circle (0.38);
\draw[-{Stealth[length=5pt]}, line width=1.8pt, obs] (6.45,2.85) ++(230:0.2) -- ++(50:0.5);
\fill[white, opacity=0.85] (0.55,0.5) circle (0.32); \fill[black!70] (0.43,0.5) -- (0.65,0.63) -- (0.65,0.37) -- cycle;
\fill[white, opacity=0.85] (1.4,0.5) circle (0.32); \fill[black!70] (1.52,0.5) -- (1.3,0.63) -- (1.3,0.37) -- cycle;
\node[fill=obs!80, text=white, rounded corners=3pt, font=\tiny\bfseries, minimum width=8mm, minimum height=6mm] at (5.6,0.5) {Brake};
\node[fill=cli, text=white, rounded corners=3pt, font=\tiny\bfseries, minimum width=8mm, minimum height=8mm] at (6.5,0.6) {Gas};
\node[fill=tC, text=black, rounded corners=3pt, font=\tiny\bfseries, minimum width=9mm, minimum height=5mm] at (6.15,1.45) {Deliver};
\node[fill=white, text=black, rounded corners=3pt, font=\tiny\bfseries, minimum width=9mm, minimum height=5mm] at (0.9,1.45) {Find me};}

% ---- side cut of the board (14 units wide) ----
\newcommand{\boardside}{%
\fill[brown!20] (-0.6,-0.35) rectangle (14.6,0);
\foreach \x/\h/\c in {1.0/0.7/black!18, 2.2/1.1/black!22, 4.6/1.9/black!25, 5.8/2.6/black!28, 7.4/2.2/black!25, 8.6/1.6/black!22, 10.9/1.0/black!18}
  \filldraw[fill=\c, draw=black!45] (\x,1.5) rectangle (\x+0.9,1.5+\h);
\filldraw[fill=tB!35, draw=tB!70!black] (12.4,1.5) rectangle (13.4,2.3);
\fill[black!70, rounded corners=1pt] (0.5,0) rectangle (1.1,0.9);
\fill[black!70, rounded corners=1pt] (12.9,0) rectangle (13.5,0.9);
\filldraw[fill=brown!45, draw=brown!70!black] (0,0.9) rectangle (14,1.3);
\fill[black!80] (0,1.3) rectangle (14,1.38);
\foreach \x in {0.3,0.9,...,13.8} \fill[black!35] (\x,1.42) circle (0.06);
\foreach \x/\c in {3.3/tA, 3.9/tA, 6.9/obs, 7.5/obs, 10.5/tC} \fill[\c] (\x,1.42) circle (0.07);
\fill[white, opacity=0.55] (0,1.38) rectangle (14,1.52);
\filldraw[fill=brd!40, draw=brd] (2.3,0.35) rectangle (4.3,0.85);
\node[font=\tiny] at (3.3,0.6) {ESP32};
\filldraw[fill=brd, draw=brd!60!black] (6.75,0.62) rectangle (7.25,0.88);
\filldraw[fill=adm!60, draw=adm!70!black] (1.8,0.6) circle (0.24);
\filldraw[fill=adm!60, draw=adm!70!black] (12.2,0.6) circle (0.24);
\filldraw[fill=black!15, draw=black!50] (9.3,0.35) rectangle (10.9,0.85);
\node[font=\tiny] at (10.1,0.6) {Power};}

\begin{document}

\thispagestyle{ckfirst}
\begin{center}
{\bfseries GEBZE TECHNICAL UNIVERSITY\par}
{Computer Engineering Department\par}
{CSE 396 Module Documentation\par}
\vspace{8pt}
{\LARGE\bfseries Çılgın Kurye\par}
\vspace{3pt}
{\large\itshape Module M2: Network and Synchronization\par}
\vspace{6pt}
\textbf{Instructor:} Erkan Zergeroğlu\hspace{1.4em}\textbf{Company:} Board\&Byte\hspace{1.4em}\textbf{Group:} Group-8\hspace{1.4em}\textbf{Date:} 6 November 2026
\end{center}
\begin{center}\footnotesize\renewcommand{\arraystretch}{1.15}
\begin{tabular}{|p{0.27\linewidth}|p{0.15\linewidth}|p{0.27\linewidth}|p{0.15\linewidth}|}\hline
\textbf{Name} & \textbf{Student No.} & \textbf{Name} & \textbf{Student No.} \\ \hline
Tuğberk Akbulut & 230104004070 & Ayşe Poyraz & 230104004085 \\ \hline
Hasan Hüseyin Özdemir & 230104004048 & Umutcan Ocak & 220104004903 \\ \hline
\end{tabular}
\end{center}
\vspace{2pt}



% ---- table and list helpers ----
\newcolumntype{L}[1]{>{\raggedright\arraybackslash}p{\dimexpr#1\linewidth-2\tabcolsep-\arrayrulewidth\relax}}
\newcolumntype{C}[1]{>{\centering\arraybackslash}p{\dimexpr#1\linewidth-2\tabcolsep-\arrayrulewidth\relax}}
\newcommand{\ckdone}{\item[$\boxtimes$]}
\newcommand{\ckopen}{\item[$\square$]}
\newcommand{\mdot}{$\bullet$}
\newcounter{msec}
\newcounter{chap}
% ---- end of helpers ----
\newcommand{\msection}[1]{\stepcounter{msec}\section{\arabic{msec}. #1}}

\msection{Purpose}

M2 connects the phones, the admin page and the LED board to the game server. It does not make game
decisions. It takes messages from outside, checks them and gives them to M1. After each tick it
takes the state from M1 and sends it to everybody. M2 also makes sure that nothing from outside can
change the game state directly. For the admin, it measures the network and it can add delay and
packet loss on purpose.

\begin{longtable}[]{|L{0.18}|L{0.82}|}
\hline Runs on & Main computer, in the same server program as M1 \\ \hline
Technology & TypeScript on Node.js; WebSocket for phones and admin; TCP and UDP for the LED board; MessagePack encoding; mDNS announcement; QR code for joining \\ \hline
Requirements & R1 R9 R10 R11 \\ \hline
\end{longtable}

\begin{figure}[htbp]\centering
\systemmap{M2}
\caption{Where M2 sits in the system. All traffic from outside passes through M2.}
\end{figure}

\msection{Requirements covered}

\begin{longtable}[]{|L{0.07}|L{0.25}|L{0.68}|}
\hline \textbf{ID} & \textbf{Requirement} & \textbf{What M2 does for it} \\ \hline \endhead
R1 & Multiplayer game & Up to 8 phones join with a QR code or a 6-digit game code, only with the browser \\ \hline
R9 & Live admin panel & Admin channel for commands, parameters, state, metrics and board data \\ \hline
R10 & Network and measurement & Ping, packet loss and frame rate for each player; added delay and loss for each player and for the board \\ \hline
R11 & Single authority & Inputs are checked and queued; every phone gets the same frozen state of each tick \\ \hline
\end{longtable}

\msection{Inputs and outputs}

\begin{longtable}[]{|L{0.08}|L{0.20}|L{0.18}|L{0.54}|}
\hline \textbf{Dir.} & \textbf{Data} & \textbf{From / to} & \textbf{Content} \\ \hline \endhead
In & Phone messages & M5 & Hello with player key, name and game code; join with team; controls; ``Find me''; ping replies \\ \hline
In & Admin messages & M6 & Login, commands, parameter changes, requests for parameters and the game log \\ \hline
In & Board messages & M7 & Earthquake events, sensor samples, board status, delay replies \\ \hline
In & State and metrics & M1 & Frozen game state and board state every tick; metrics twice per second \\ \hline
Out & To M1 & M1 & Checked controls, joins, ``Find me'', admin commands, board events; network metrics \\ \hline
Out & To phones & M5 & Game state 30 times per second; join replies; ping; frame-rate limit \\ \hline
Out & To admin & M6 & State 10 times per second, metrics, replies, board data \\ \hline
Out & To board & M7 & Board state 10 to 30 times per second; config; vibrate commands \\ \hline
\end{longtable}

\msection{Main components}

\begin{longtable}[]{|L{0.20}|L{0.52}|L{0.28}|}
\hline \textbf{Component} & \textbf{What it does} & \textbf{Related modules} \\ \hline \endhead
Web server and join & Serves the phone and admin pages, the QR code and the game code; announces the server on the local network & M5, M6, M7 \\ \hline
Sessions & One session for each phone and the admin; player key for reconnecting within 30 s; one admin at a time & M1, M5, M6 \\ \hline
Codec and checker & MessagePack encode and decode; checks type, fields and value ranges; a bad message gets an error reply & all \\ \hline
Input queue & Holds checked inputs until M1 reads them at the start of the next tick & M1 \\ \hline
State broadcaster & Encodes the frozen state once and sends the same bytes to every phone & M1, M5, M6, M7 \\ \hline
Delay and loss emulator & An up queue and a down queue for each player and for the board, set by the admin & M6 \\ \hline
Ping and metrics & Ping twice per second, clock offset for the phone, ping, loss and frame-rate metrics & M1, M5, M6 \\ \hline
Board link & Server side of the board connection: one TCP channel and one UDP stream & M7 \\ \hline
\end{longtable}

\msection{How it works}

\textbf{Which link uses which protocol.} This choice is important, so we made it message by message.
A web browser can only open a WebSocket for two-way talk, so phones and the admin page use
WebSocket. WebSocket runs on TCP, so every message arrives and the order is kept. The LED board is
not a browser, so it can use both TCP and UDP. On TCP it gets the messages that must never be lost:
the config (player colours, brightness, current limit) and the vibrate commands. The board also
sends earthquake events and its status report on TCP. If one of these is lost, the game goes wrong,
so TCP sends it again. The fast and repeated data goes on UDP: the board state, which comes 10 to 30
times per second, the sensor samples for the seismic graph, and the small delay probes. Here only
the newest message is useful. If one UDP packet is lost, the next one comes 33 to 100 ms later and
replaces it. TCP would wait for the lost packet and the LEDs would stop for a moment, so UDP is the
better choice for this data. Each board state fits in one UDP packet (under 1 KB), so a packet is
never split.

\begin{figure}[htbp]\centering
\begin{tikzpicture}[ck]
\fill[tcpc!8, rounded corners=4pt] (2.0,0.95) rectangle (11.0,-0.95);
\fill[tcpc!8, rounded corners=4pt] (2.0,-1.35) rectangle (11.0,-3.25);
\fill[tcpc!14, rounded corners=4pt] (2.0,-3.65) rectangle (11.0,-5.55);
\fill[udpc!12, rounded corners=4pt] (2.0,-5.95) rectangle (11.0,-8.25);
\node[font=\scriptsize\bfseries, text=tcpc] at (6.5,0.72) {WebSocket on TCP};
\node[font=\scriptsize\bfseries, text=tcpc] at (6.5,-1.58) {WebSocket on TCP};
\node[font=\scriptsize\bfseries, text=tcpc] at (6.5,-3.88) {TCP: must arrive};
\node[font=\scriptsize\bfseries, text=udpc] at (6.5,-6.18) {UDP: newest wins};
\node[sv, minimum width=24mm, minimum height=91mm, line width=2pt] (m2) at (0.6,-3.65) {\textbf{M2}\\on the main\\computer};
\node[cl, minimum width=22mm, minimum height=16mm] (ph) at (12.4,0) {\textbf{M5}\\Phones};
\node[ad, minimum width=22mm, minimum height=16mm] (ad) at (12.4,-2.3) {\textbf{M6}\\Admin page};
\node[bd, minimum width=22mm, minimum height=39mm] (bo) at (12.4,-5.9) {\textbf{M7}\\LED board};
\draw[tcpa,->] (2.0,0.2) -- node[al]{game state, 30 per s} (11.0,0.2);
\draw[tcpa,<-] (2.0,-0.4) -- node[al]{controls, ``Find me''} (11.0,-0.4);
\draw[tcpa,->] (2.0,-2.1) -- node[al]{state, metrics, board data} (11.0,-2.1);
\draw[tcpa,<-] (2.0,-2.7) -- node[al]{commands, parameters} (11.0,-2.7);
\draw[tcpa,->] (2.0,-4.4) -- node[al]{config, vibrate} (11.0,-4.4);
\draw[tcpa,<-] (2.0,-5.0) -- node[al]{earthquake event, status} (11.0,-5.0);
\draw[udpa,->] (2.0,-6.7) -- node[al]{board state, 10 to 30 per s} (11.0,-6.7);
\draw[udpa,<-] (2.0,-7.25) -- node[al]{sensor samples} (11.0,-7.25);
\draw[udpa,<->] (2.0,-7.8) -- node[al]{delay probe} (11.0,-7.8);
\end{tikzpicture}
\caption{Every link and every message type. Blue solid arrows use TCP, orange dashed arrows use UDP.}
\end{figure}

\textbf{Joining.} At the start the admin screen shows a QR code and a 6-digit game code. A player
scans the QR code with the phone camera. The phone opens the game page from the main computer, so
there is nothing to install. A player without a camera can open the short local address and type
the code. The phone then chooses a team. The first time, the browser saves a random player key. If
the page reloads or Wi-Fi drops, the phone sends the same key and gets the same motorcycle back
within 30 seconds. The LED board finds the server by itself, because the server announces its name
and port on the local network with mDNS.

\begin{figure}[htbp]\centering
\begin{tikzpicture}[ck]
\fill[black!80, rounded corners=4pt] (-0.2,-0.2) rectangle (6.2,4.2);
\fill[white] (0,0) rectangle (6,4);
\fill[srv] (0,3.55) rectangle (6,4);
\node[text=white, font=\scriptsize\bfseries, anchor=west] at (0.15,3.78) {Çılgın Kurye, admin};
\node[font=\scriptsize\bfseries] at (1.75,3.25) {Scan to join};
\begin{scope}[shift={(0.55,0.55)}]
\pgfmathsetseed{7}
\foreach \x in {0,...,20}{\foreach \y in {0,...,20}{%
  \pgfmathparse{rnd>0.52?1:0}\ifnum\pgfmathresult=1 \fill[black] (\x*0.11,\y*0.11) rectangle ++(0.11,0.11);\fi}}
\foreach \x/\y in {0/0,0/14,14/14}{\fill[white] (\x*0.11-0.02,\y*0.11-0.02) rectangle ++(0.81,0.81);
  \fill[black] (\x*0.11,\y*0.11) rectangle ++(0.77,0.77); \fill[white] (\x*0.11+0.11,\y*0.11+0.11) rectangle ++(0.55,0.55); \fill[black] (\x*0.11+0.22,\y*0.11+0.22) rectangle ++(0.33,0.33);}
\end{scope}
\node[font=\scriptsize, anchor=west] at (3.3,2.9) {or open};
\node[font=\scriptsize\ttfamily, anchor=west] at (3.3,2.5) {kurye.local};
\node[font=\scriptsize, anchor=west] at (3.3,1.9) {Game code};
\node[draw=srv, fill=srvL, font=\large\bfseries, rounded corners=2pt, anchor=west] at (3.3,1.3) {482 913};
\node[font=\tiny, anchor=west, text=black!60] at (3.3,0.45) {Players: 3 of 8};
\begin{scope}[shift={(8.0,-0.3)}]
\fill[black!85, rounded corners=7pt] (-0.15,-0.15) rectangle (2.55,4.75);
\fill[white] (0,0) rectangle (2.4,4.6);
\node[font=\scriptsize\bfseries] at (1.2,4.25) {Join the game};
\node[draw=black!40, minimum width=20mm, font=\tiny, anchor=west, text=black!55] at (0.2,3.7) {Name: Deniz};
\node[draw=black!40, minimum width=20mm, font=\tiny, anchor=west, text=black!55] at (0.2,3.2) {Code: 482 913};
\node[font=\tiny\bfseries] at (1.2,2.75) {Choose a team};
\node[fill=tA, text=white, font=\tiny\bfseries, minimum width=9mm, minimum height=6mm, rounded corners=2pt] at (0.65,2.25) {Blue};
\node[fill=tB, text=white, font=\tiny\bfseries, minimum width=9mm, minimum height=6mm, rounded corners=2pt] at (1.75,2.25) {Green};
\node[fill=tC, text=black, font=\tiny\bfseries, minimum width=9mm, minimum height=6mm, rounded corners=2pt] at (0.65,1.55) {Yellow};
\node[fill=tD, text=white, font=\tiny\bfseries, minimum width=9mm, minimum height=6mm, rounded corners=2pt] at (1.75,1.55) {Purple};
\node[fill=cli, text=white, font=\tiny\bfseries, minimum width=20mm, minimum height=6mm, rounded corners=2pt] at (1.2,0.6) {Ready};
\end{scope}
\draw[inna,->, line width=1.2pt] (6.3,1.9) -- node[al]{scan} (7.75,1.9);
\end{tikzpicture}
\caption{Joining: the QR code and game code on the admin screen, and the join page on a phone.}
\end{figure}

\textbf{Keeping the game state safe.} Many phones send messages at the same time, but only M1 may
change the game state, and only inside a tick. So M2 never writes to the state. A phone can only
send controls; it can never send a position or a score. M2 decodes each message, checks its type
and fields, and cuts every value into its range (for example steer is kept between −1 and 1). Each
session can only control its own player, so one phone cannot move another motorcycle. Every message
has an order number. A message with an old or repeated number is dropped. Good inputs wait in a
queue, and M1 reads the queue at the start of the next tick. At the end of the tick M1 gives a
frozen copy of the state. M2 encodes it once with MessagePack and sends the same bytes to all phones,
so all players see exactly the same tick. The full game state is under 4 KB.

\begin{figure}[htbp]\centering
\begin{tikzpicture}[ck, node distance=4mm]
\node[cl, minimum width=16mm] (p) {Phone\\message};
\node[sv, minimum width=16mm, right=of p] (c) {Decode\\and check};
\node[sv, minimum width=16mm, right=of c] (q) {Input\\queue};
\node[sv, minimum width=16mm, right=of q, fill=srv!35, line width=1.6pt] (t) {M1 tick};
\node[sv, minimum width=16mm, right=of t] (f) {Frozen\\state};
\node[sv, minimum width=16mm, right=of f] (e) {Encode\\once};
\node[cl, minimum width=16mm, right=of e] (a) {All\\phones};
\node[nb, draw=obs, fill=obs!10, below=8mm of c] (x) {Dropped,\\error reply};
\foreach \u/\v in {p/c,c/q,q/t,t/f,f/e} \draw[inna,->] (\u) -- (\v);
\draw[tcpa,->] (e) -- (a);
\draw[inna,->, obs] (c) -- node[al]{bad} (x);
\end{tikzpicture}
\caption{The path of an input. Only M1 writes the state, and everybody gets the same bytes.}
\end{figure}

\textbf{Ping, delay and loss.} Twice per second M2 sends a ping with the server time. The phone
answers with its own time, its frame rate and the number of the last state it got. From this M2
finds the round-trip time and the clock offset of the phone. Gaps in the order numbers give the
packet loss. These values go to the admin graphs. For tests the admin can make the network worse
for one player. Each player has an up queue and a down queue. A queue drops a message with the loss
chance that the admin set, and holds the other messages for half of the added ping. The same tool
works on the board link. The admin can also set a frame-rate limit, which M2 sends to the phone.

\msection{Testing and demonstration}

M2 is tested with a fake M1 that sends a moving test state, with up to 8 scripted bot phones, and
with the board emulator of M7. The tests check the message codec, bad and broken input, the load
with 8 phones, how exact the added delay and loss are, reconnecting within 30 seconds, and the board
link with 20\% packet loss. In the demonstration phones join with the QR code or the game code. Then
the admin changes one player's ping and loss live, and we watch the change on that phone and in the
graphs.

\end{document}
